
\subsection{Round five: method}
\label{sec:r5-method}
\paragraph{Question N Q1.6.} On designs with a known optimal clustering, how far are the VPR packer and the Norns HDL-aware packer from it, and how does the gap grow with size? The idea follows Cong, Romesis and Xie~\cite{Cong03tpeko}, who built placement problems with known optima to measure how far placers are from them.
\paragraph{What ``optimal'' means here.} A \emph{clustering} optimum by construction, not a proven placement optimum. A block is 10 six-input functions, each feeding a flip-flop (10 logic elements of the \code{k6\_frac\_N10} architecture, so exactly one CLB), with dense internal connections and at most two external inputs per function. The CLB count $B$ is an exact lower bound and is reached by the block partition. The block partition's cut is exact for that partition but is an upper bound on the minimum cut: at the netlist level a local search of improving pair swaps found a slack of 0.0\% to 6.8\% (median 0.0\%) on the 72 netlist designs. The ``known'' arm hands the block partition to the same placer and router, so its wirelength, channel width and critical-path delay are what VPR achieves from a good clustering, not a bound; ratios below one are possible.
\paragraph{Designs.} Nine sizes ($B=4$ to $1024$), four topologies (chain, mesh, tree, random sparse), two styles (\code{dp}: a datapath slice whose HDL has the vector \code{r[9:0]}; \code{ctl}: scalar registers, random logic with no vector), generated deterministically from the seed 20261026. Two levels: NL, the netlist written directly as BLIF (no Odin III, no ABC), and HDL, the same specifications as Verilog through Odin III. Arms: \code{stock} VPR; \code{norns} (HDL level only: the round-four headline packer, affinity source \code{word}, $\beta=4$, which needs provenance); \code{known} (the Norns binary with the block membership as affinity groups, $\beta=100$). Three VPR seeds each. Synthesis keeps the construction in 63 of 72 HDL designs (original label rule: 32); the netlist check passes on 72 of 72. Excluded by the rule fixed in advance: r5\_chain\_ctl\_1024, r5\_chain\_ctl\_512, r5\_chain\_ctl\_64, r5\_mesh\_ctl\_1024, r5\_mesh\_ctl\_256, r5\_mesh\_ctl\_512, r5\_tree\_ctl\_1024, r5\_tree\_ctl\_32, r5\_tree\_ctl\_512.
Rules, thresholds and statuses are in the protocol (\code{docs/plan/round5.md}, SHA \code{7886f49}) and four dated entries; the paper repeats none of them.
